% ============================================================
%  CHAPTER 7 — Discussion
% ============================================================

\chapter{Discussion}
\label{ch:discussion}

% ============================================================
\section{Validity of Results}
% ============================================================

\subsection{Embedded Subsystem}

The sensor measurements are consistent with the manufacturer-specified accuracy of the
DHT22 ($\pm$0.5\,°C, $\pm$2--5\,\% RH) and within the expected operational range of
the MAX30102 for resting-state adult subjects.
The BPM standard deviation of 2.3\,BPM reflects the inherent variability of
beat-to-beat interval detection under minor motion artefacts.
The SpO\textsubscript{2} accuracy of $\pm$1\,\% is within the $\pm$2\,\% tolerance
generally accepted.

% -------------------------------------------------------
% TODO (Backend Team): discuss the validity of your API
% test results (is the response time consistent? are there
% edge cases that were not covered? is the alert logic
% triggering correctly?).
% -------------------------------------------------------
\begin{center}
	\textit{[Backend results validity discussion to be added.]}
\end{center}

% -------------------------------------------------------
% TODO (Web / AI Teams): discuss validity of your results.
% -------------------------------------------------------
\begin{center}
	\textit{[Web dashboard and AI results validity discussion to be added.]}
\end{center}

% ============================================================
\section{Limitations}
% ============================================================

\subsection{Embedded Subsystem}

\begin{description}
	\item[No local offline buffer.]
	      Records generated during network outages are discarded.

	\item[Plaintext NVS storage.]
	      Firebase credentials (API key, email, password) are stored in NVS without
	      encryption.
	      The ESP32 NVS partition supports native encryption via an NVS key partition;
	      this should be enabled before any production deployment.

	\item[Implicit shared-memory atomicity.]
	      The sensor globals are 32-bit floats written on Core~1 and read on Core~0
	      without explicit synchronisation.
	      A FreeRTOS mutex or atomic qualifier is the correct long-term solution.

	\item[Single patient per device.]
	      The current design maps one device to exactly one patient and one room.
	      Multi-patient rooms would require either multiple devices or a structural change
	      to the data model.
\end{description}

% -------------------------------------------------------
% TODO (Backend Team): list and discuss backend limitations
% (e.g. no rate limiting, single-region deployment, no
% caching layer, alert delivery latency, etc.).
% -------------------------------------------------------
\subsection{Backend Subsystem}
\begin{center}
	\textit{[To be completed by the backend team.]}
\end{center}

% -------------------------------------------------------
% TODO (Web Dashboard Team): list and discuss dashboard
% limitations (e.g. no offline mode, supported browsers,
% chart rendering performance with large datasets, etc.).
% -------------------------------------------------------
\subsection{Web Dashboard Subsystem}
\begin{center}
	\textit{[To be completed by the web dashboard team.]}
\end{center}

% -------------------------------------------------------
% TODO (AI Team): list and discuss AI limitations
% (e.g. limited training data, class imbalance, model
% latency budget, false positive rate, dependency on
% data quality from the embedded node, etc.).
% -------------------------------------------------------
\subsection{AI Subsystem}
\begin{center}
	\textit{[To be completed by the AI team.]}
\end{center}

% ============================================================
\section{ESP32 Constraints Impact}
% ============================================================

The most significant platform constraint encountered was the non-thread-safety of
the mbedTLS context, which imposed the non-obvious architectural requirement of
deferring Firebase initialisation until the task is running on Core~0.

The 8\,KiB task stack was found to be adequate for the TLS handshake and async send
operations, with an estimated 1.5\,KiB headroom.

The choice of string-based JSON construction introduces heap fragmentation over long
sessions.

% ============================================================
\section{Robustness}
% ============================================================

The unified \texttt{wifiConnect()} function using \texttt{WiFi.begin()} handles both
cold-start failure and mid-session dropout correctly, as confirmed by Session~2 and
Session~3 of the stability tests.

The delta-threshold filter provides robustness against spurious sensor noise:
fluctuations smaller than 0.3\,°C, 0.5\,\% RH, or 2\,\% SpO\textsubscript{2} do
not generate database writes, reducing both network load and the volume of
low-information records in the database.

% ============================================================
\section{Identified Improvements}
% ============================================================

Priority improvements for future versions of the embedded subsystem:

\begin{enumerate}
	\item \textbf{Local offline buffer.}
	      A LittleFS-backed FIFO queue for records generated during network outages,
	      with automatic retry on reconnection.
	      This would eliminate the record loss observed in Sessions~2 and~3.

	\item \textbf{NVS encryption.}
	      Enable the ESP32 native NVS encryption partition to protect credentials at rest.

	\item \textbf{Thread-safe sensor globals.}
	      Replace raw float globals with a mutex-protected struct or a FreeRTOS queue
	      for safe inter-core communication.

	\item \textbf{Wi-Fi modem sleep.}
	      Implement modem sleep between upload intervals to reduce average current usage for future battery-powered deployments.
\end{enumerate}

% -------------------------------------------------------
% TODO (Backend / Web / AI Teams): add your own improvement
% lists here.
% -------------------------------------------------------
\begin{center}
	\textit{[Improvements for the backend, web dashboard, and AI subsystems to be added.]}
\end{center}
